\section{What the final theorem does not use}\label{sec:offpath}

Nearly every file of the archive imports nearly every other, so the import graph says nothing
about what depends on what. The statements of this section are instead computed from the
compiled archive: the \emph{cone} of \lean{riemannHypothesis} is the set of all declarations
reachable from it through the type and the value of every declaration reached
(\cref{sec:cone}). It has \conesize{} declarations, of which \conemathlib{} belong to
Mathlib. The other four theorems of \lean{Main.lean} are derived from the same statement for
all characters, by short arguments of their own: the cone of all five has \conefive{}
declarations, twenty-four more.

\begin{table}[ht]
\centering
\begin{tabular}{@{}lrrr@{}}
\toprule
 & in the archive & in the cone & share\\
\midrule
Theorems & \projtheorems & \coneprojtheorems & 69\%\\
Definitions & \projdefs & \coneprojdefs & 54\%\\
Modules that declare something & \projmodules & \coneprojmodules & 84\%\\
\bottomrule
\end{tabular}
\caption{The archive and the cone of \texttt{riemannHypothesis}. A module is counted in the
cone if it declares at least one declaration of the cone.}\label{tab:cone}
\end{table}

So about a third of the archive's theorems, 494{,}644 of them, are not used by the final
theorem. They are compiled and checked by the kernel like the rest
(\cref{sec:verification}); they are simply not on the path. By top-level directory the unused
modules number 22{,}037 of 141{,}928 in \lean{Development}, 3{,}224 of 18{,}254 in
\lean{Formalism}, 1{,}814 of 8{,}507 in \lean{GroundTheory}, 360 of 3{,}845 in
\lean{ReciprocityTheory}, 140 of 2{,}027 in \lean{CompanionTheory} and 106 of 449 in
\lean{PositivityTheory}.

What is unused includes earlier and parallel forms of steps that are on the path. Each of the
following was confirmed outside the cone by its full name.
\begin{itemize}
\item \emph{Conditional reductions.} Several statements of the form ``if such data exist at
  every scale, then the Riemann Hypothesis holds'':
  \begin{center}\small
  \begin{tabular}{@{}ll@{}}
  \lean{RiemannHypothesisOfLowerBounds} & \lean{RiemannHypothesisOfHalfLinks}\\
  \lean{WeightedSieveImpliesRH} & \lean{WeightedSieveImpliesRHThreeParameter}\\
  \lean{GrowthBoundGivesRiemannHypothesis} & \lean{RiemannHypothesisOfBalanceAndGrowthData}
  \end{tabular}
  \end{center}
  One more of this kind, \lean{ZetaQuarterPowerCorollary}, takes the number of lifts
  $m\approx x^{1/4}$ and concludes only that $\zeta(s)\neq0$ for $\re s>\tfrac34$.
\item \emph{A first form of the analytic half.} The identity of \cref{st:balance} and the
  reduction of \cref{sec:reduction} exist a second time, as \lean{AskMainBalance} (proved as
  \lean{colour_balance}) and \lean{GrowthInequalityForm}, for hosts ramified at most at $2$
  and two odd primes. The path uses the versions that allow a set $S$ of primes in place of
  $2$ and the first odd prime (\lean{ch_balance}, \lean{ch_core}).
\item \emph{Two further copies of the gain formula} of \cref{sec:adm}, in other namespaces.
\item \emph{Other forms of single steps:} the jet statement \lean{thmBGauss}; a second form
  of the mean bound, \lean{MeanBoundSecondForm}; the determinant statements \lean{DetGain}
  (if $g_p\le v_p(\Delta)$ on a set of primes then $\sum g_p\log p\le\log\abs\Delta$),
  \lean{DetChain} and \lean{DetBlocks}; a Bessel period
  statement \lean{BesselPeriodGGP}; the congruence-number statements for $\GSp_4$
  (\lean{NumberSp4}) and the first form for $\U(3)$ (\lean{CongruenceNumberRankThree});
  and one of the derivations of the lower bound for the scalar of the lifts,
  \lean{lamLow_of_classicalTheorems}.
\end{itemize}

In the other direction, every theorem declared in the file that collects the results taken from
the literature (\cref{sec:literature}), \hubtheorems{} in all with the auxiliary ones, is in the cone. None of that
library is decoration.

The archive declares no axiom of its own, on the path or off it: among its 1{,}945{,}519
declarations there are theorems, definitions, inductive types with their constructors and
recursors, and two opaque constants, and nothing else.
